\section*{Chapter \ref{ch:b1:lipids} --- Lipids}

\begin{solution}{exo:b1:lipids:1}
A \omterm{def:b1:lipids:fattyacid}{lipid} is a biological molecule insoluble in water and soluble in
non-polar solvents. \omterm{def:b1:lipids:triglyceride}{Fats}, \omterm{def:b1:lipids:amphiphilic}{phospholipids}, \omterm{def:b1:lipids:amphiphilic}{sterols} and carotenoids share
no common skeleton, only this behaviour toward water, which is what
gives them their common roles (stores, films, membranes).
\end{solution}

\begin{solution}{exo:b1:lipids:2}
C18:2 $\Delta^{9,12}$. The last double bond starts at carbon 12 from
the carboxyl, i.e.\ at carbon 6 from the methyl end: omega-6.
\end{solution}

\begin{solution}{exo:b1:lipids:3}
A \omterm{def:b1:lipids:triglyceride}{triglyceride} has no polar head: entirely hydrophobic, it is
excluded from water as a bulk phase, a droplet. A \omterm{def:b1:lipids:amphiphilic}{phospholipid} is
\omterm{def:b1:lipids:amphiphilic}{amphiphilic}: its polar head must stay in water while its tails must
leave it, which only a sheet two molecules thick can satisfy.
\end{solution}

\begin{solution}{exo:b1:lipids:4}
About \qty{41}{\celsius}, \qty{-5}{\celsius} and \qty{20}{\celsius}
(the \omterm{def:b1:lipids:amphiphilic}{cholesterol} curve has no sharp transition; half its rise is near
\qty{20}{\celsius}). A bacterium at \qty{10}{\celsius} needs the
unsaturated composition, fluid well below its growth temperature.
\end{solution}

\begin{solution}{exo:b1:lipids:5}
$15\,000\times 38 = \qty{570}{MJ}$; $570/8 = 71$ days. Glycogen: dry
$570\,000/17 = \qty{33.5}{kg}$, hydrated \qty{134}{kg}.
\end{solution}

\begin{solution}{exo:b1:lipids:6}
C22:0 $>$ C18:0 $>$ C14:0 $>$ C18:1 $>$ C18:3. Among saturated acids,
longer chains make more \omterm{def:b1:water-small-molecules:bonds}{van der Waals} contacts; one \emph{cis} double
bond removes more packing than four extra carbons add (C18:1 melts
below C14:0); each further bond lowers it again.
\end{solution}

\begin{solution}{exo:b1:lipids:7}
$2\times\qty{140}{\micro m^2}/\qty{0.6}{nm^2} = 2\times 1.4\times
10^{-10}/6\times 10^{-19} = \num{4.7e8}$ molecules; mass
$\num{4.7e8}\times 750/\num{6e23} = \qty{5.8e-13}{g} = \qty{0.58}{pg}$.
\end{solution}

\begin{solution}{exo:b1:lipids:8}
Volume \qty{1.11}{mL} $= \qty{1.11e-6}{m^3}$; surface of spheres
$= 6V/d = 6\times 1.11\times 10^{-6}/10^{-6} = \qty{6.7}{m^2}$. One
drop of that volume: $d = \qty{1.28}{cm}$, surface \qty{5.2}{cm^2}: the
emulsion has thirteen thousand times more interface. Lipase is
water-soluble and acts only at the interface; the rate is proportional
to it.
\end{solution}

\begin{solution}{exo:b1:lipids:9}
Above $T_m$ its rigid rings hinder the motion of the neighbouring
chains and reduce \omterm{def:b1:lipids:fluidity}{fluidity}; below $T_m$ its bulk prevents the chains
from packing into the ordered gel and prevents freezing. With no
cooperative packing possible, there is no temperature at which the
whole bilayer changes state at once: the transition is spread out.
\end{solution}

\begin{solution}{exo:b1:lipids:10}
All-saturated chains: at \qty{25}{\celsius} the membranes are already
close to their transition and stiff; at \qty{5}{\celsius} they gel.
Gelled membranes let proteins stop working and, on rewarming or
mechanical stress, crack and leak: the \omterm{def:b1:cell-unit-of-life:cell}{cells} lose their contents and
the \omterm{def:b1:body-plans-tissues:tissue}{tissue} dies (chilling injury). The wild type desaturates its
\omterm{def:b1:lipids:fattyacid}{lipids} in autumn and stays fluid at \qty{5}{\celsius}.
\end{solution}

\begin{solution}{exo:b1:lipids:11}
\qty{1.07}{kg} of water per kilogram of \omterm{def:b1:lipids:triglyceride}{fat}. Oxygen: \qty{2000}{L};
at \qty{5}{\%} extraction of \qty{21}{\%} oxygen, air needed
$2000/(0.21\times 0.05) = \qty{190}{m^3}$; water lost at
\qty{30}{g/m^3}: \qty{5.7}{kg}. The camel loses five times more water
breathing than the \omterm{def:b1:lipids:triglyceride}{fat} yields: the hump is an energy store, not a water
store.
\end{solution}

\begin{solution}{exo:b1:lipids:12}
Given a bilayer, the \omterm{prop:b1:water-small-molecules:properties}{hydrophobic effect} makes it reseal, close into
vesicles and accept new \omterm{def:b1:lipids:fattyacid}{lipids} with no energy input: assembly is
spontaneous. But the \omterm{def:b1:cell-unit-of-life:cell}{cell} makes its \omterm{def:b1:lipids:fattyacid}{lipids} with enzymes sitting in an
existing membrane (the ER), which insert them in place; a new membrane
grows by expansion of an old one and is partitioned at division, so
that every membrane of every \omterm{def:b1:cell-unit-of-life:cell}{cell} descends from the membranes of the
previous \omterm{def:b1:cell-unit-of-life:cell}{cell}. Self-assembly explains the physics of the sheet, not its
origin in the \omterm{def:b1:cell-unit-of-life:cell}{cell}.
\end{solution}

\begin{solution}{pb:b1:lipids:1}
\textbf{1.} C16:0 (0), C18:1 $\Delta^9$ (1), C22:6 $\Delta^{4,7,10,13,16,19}$
(6).
\textbf{2.} Warm: $0.40\times 0 + 0.45\times 1 + 0.15\times 6 = 1.35$;
cold: $0 + 0.40 + 2.10 = 2.50$ double bonds per chain.
\textbf{3.} Each \emph{cis} bond kinks the chain and prevents close
packing with neighbours; fewer \omterm{def:b1:water-small-molecules:bonds}{van der Waals} contacts, less heat needed
to disorder the chains, lower $T_m$.
\textbf{4.} $1.15$ more bonds per chain: about \qty{45}{\celsius}
lower.
\textbf{5.} In a gel the \omterm{def:b1:lipids:fattyacid}{lipids} cannot move: carriers cannot change
conformation, pumps stop, \omterm{def:b1:membranes-transport:transporters}{channels} stick; ion gradients decay, nerve
conduction fails, and the fish is paralysed and dies of the cold that
a cold-acclimated fish tolerates.
\textbf{6.} Desaturases, which introduce double bonds into fatty acyl
chains; they are \omterm{def:b1:membranes-transport:proteins}{integral proteins} of the \omterm{def:b1:eukaryotic-cell:endomembrane}{endoplasmic reticulum}, where
the \omterm{def:b1:lipids:fattyacid}{lipids} are made.
\textbf{7.} In both it \omterm{def:b1:water-small-molecules:buffer}{buffers} \omterm{def:b1:lipids:fluidity}{fluidity}: stiffening the very
unsaturated cold membrane above its low $T_m$, and keeping the warm
membrane from gelling on a cool day.
\textbf{8.} $30\,000\times 38 = \qty{1140}{MJ}$.
\textbf{9.} $1140/60 = 19$ days.
\textbf{10.} Dry: $\num{1.14e6}/17 = \qty{67}{kg}$; hydrated:
\qty{268}{kg}.
\textbf{11.} $268 - 30 = \qty{238}{kg}$, nearly half the camel's mass.
\textbf{12.} A layer of \omterm{def:b1:lipids:triglyceride}{fat} under the whole skin would insulate the
body and prevent it from shedding heat in the desert; concentrating
the \omterm{def:b1:lipids:triglyceride}{fat} in one hump \omterm{def:b1:flowering-plant-organization:organs}{leaves} the rest of the skin free to lose heat.
\textbf{13.} $1.07\times 30 = \qty{32}{kg}$ of water.
\textbf{14.} $2.0\times 30\,000 = \qty{60000}{L} = \qty{60}{m^3}$ of
oxygen.
\textbf{15.} $60/(0.21\times 0.05) = \qty{5700}{m^3}$ of air.
\textbf{16.} $5700\times 30 = \qty{171}{kg}$ of water.
\textbf{17.} Lost five times what is gained: oxidising the hump costs
water. The hump is a store of energy that lets the camel go without
food; its water economy comes from its kidneys, its tolerance of
dehydration and its ability to let its temperature rise.
\textbf{18.} $500\times 3.5\times 6 = \qty{10.5}{MJ}$ stored as heat and released at night; evaporating it away would have cost $10\,500/2.4 = \qty{4.4}{kg}$ of water a day.
\textbf{19.} A membrane can stretch only a few percent, so doubling the
volume of a red \omterm{def:b1:cell-unit-of-life:cell}{cell} requires a large excess of membrane folded into
the resting biconcave shape, and a \omterm{def:b1:eukaryotic-cell:cytoskeleton}{cytoskeleton} that lets it unfold
without tearing: the camel's red \omterm{def:b1:cell-unit-of-life:cell}{cells} are oval, with more membrane per
volume than a human's.
\textbf{20.} $2\times 80\times 10^{-12}/0.65\times 10^{-18} = \num{2.5e8}$.
\textbf{21.} $\num{2.5e8}\times\num{8e12}\times 40 = \num{7.9e22}$.
\textbf{22.} $\num{7.9e22}\times 750/\num{6e23} = \qty{99}{g}$.
\textbf{23.} Lateral diffusion keeps the head in water and the tails in
\omterm{def:b1:lipids:triglyceride}{oil} at every step, costing nothing; flipping requires the polar head to
cross the hydrophobic core, a large energy barrier that thermal
agitation surmounts once a week.
\textbf{24.} A difference between leaflets decays only as fast as
\omterm{def:b1:lipids:fattyacid}{lipids} flip; at once a week per molecule, an asymmetry set up by
enzymes that move specific \omterm{def:b1:lipids:fattyacid}{lipids} across (flippases) persists
indefinitely against spontaneous flipping.
\textbf{25.} About \qty{240}{kg} saved --- nearly half its body mass
--- by carrying \qty{30}{kg} of \omterm{def:b1:lipids:triglyceride}{fat} instead of \qty{270}{kg} of hydrated
glycogen; and the hump is not a water store: oxidising it loses five
times more water in the breath than it produces.
\end{solution}
